\providecommand{\Sp}{\mathrm{Sp}}
\providecommand{\Mp}{\mathrm{Mp}}
\providecommand{\Aut}{\mathrm{Aut}}
\providecommand{\Z}{\mathbb{Z}}
\providecommand{\C}{\mathbb{C}}
\providecommand{\cO}{\mathcal{O}}
\providecommand{\DT}{\mathrm{DT}}
\providecommand{\red}{\mathrm{red}}
\providecommand{\kBKM}{\kappa_{\mathrm{BKM}}}
\providecommand{\kch}{\kappa_{\mathrm{ch}}}
\providecommand{\kcat}{\kappa_{\mathrm{cat}}}
\providecommand{\kfib}{\kappa_{\mathrm{fiber}}}
\providecommand{\wt}{\operatorname{wt}}

\section{Order-$5$ Automorphic Weight and Metaplectic Roots}
\label{sec:lorgat-N5-metaplectic}

\subsection*{Three normalisations}

The primitive CHL denominator ladder over the \(K3\times E\) slice has
exactly five orders:
\[
  N\in\{1,2,3,4,6\},\qquad
  c_N(0)=(10,8,6,4,2),\qquad
  \kBKM(\Phi_N)=c_N(0)/2=(5,4,3,2,1).
\]
These are the Gritsenko--Nikulin \(\Delta\)-tower constants in the
primitive BKM-denominator convention.  The physics square convention
records the squared Siegel form; at \(N=1\),
\[
  \Delta_5=\operatorname{Bor}(\phi_{0,1}),\qquad
  \Phi_{10}=\Delta_5^2=\operatorname{Bor}(2\phi_{0,1}).
\]
Thus \(2\phi_{0,1}\) doubles the divisor multiplicities and lands on
the square.  The primitive BKM denominator uses \(\phi_{0,1}\).

The order-indexed Nikulin/Gritsenko boundary ladder is indexed by
cyclic symplectic K3 automorphisms:
\[
  N\in\{1,2,3,4,5,6,7,8\},\qquad
  c_N(0)=(10,8,6,4,4,2,2,2),
\]
hence the non-CHL boundary constants are
\[
  (c_5(0),c_7(0),c_8(0))=(4,2,2),\qquad
  (\kBKM(\Phi_5),\kBKM(\Phi_7),\kBKM(\Phi_8))=(2,1,1).
\]
For \(N=5\), the \(M_{24}\) frame shape \(1^4 5^4\) gives
\[
  T_{H^*}(g_5)=8,\qquad A_5=2,\qquad
  c_5(0)=T_{H^*}(g_5)-2A_5=4,\qquad
  \wt(\Phi_5^{\mathrm{int}})=2.
\]
The scalar \(1\) used in weight-\(1/2\) metaplectic inputs belongs to a
separately specified half-integral Jacobi or vector-valued input.  It
is never the order-indexed \(g_5\)-twined coefficient.

The Gritsenko--Cl\'ery diagonal-divisor catalogue is a third object,
indexed by paramodular triples:
\[
\begin{gathered}
 (t,N;k)=
 (1,1;5),\ (2,1;2),\ (1,2;3),\ (3,1;1),\\
 (1,3;2),\ (4,1;\tfrac12),\ (1,4;\tfrac32),\ (2,2;1).
\end{gathered}
\]
Its row weights are
\[
  (5,2,3,1,2,\tfrac12,\tfrac32,1),
\]
and its twice-weight tuple is
\[
  (10,4,6,2,4,1,3,2).
\]
This tuple belongs to the triple-indexed dd-modular catalogue.  It is
separate from the primitive CHL tuple
\((10,8,6,4,2)\) and from the order-indexed Nikulin/Gritsenko tuple
\((10,8,6,4,4,2,2,2)\).

\subsection*{The order-$5$ automorphic row}

Let \(\Phi_5^{\mathrm{int}}\) denote the order-indexed Gritsenko
paramodular form attached to the symplectic order-$5$ K3 twining.  Its
normalisation is
\[
  c_5(0)=4,\qquad
  \kBKM(\Phi_5^{\mathrm{int}})=2,\qquad
  \wt(\Phi_5^{\mathrm{int}})=2.
\]
Mukai fixed-point data, the \(M_{24}\) frame shape \(1^4 5^4\), and the
twined K3 elliptic genus give the same coefficient.  This proves the
automorphic weight statement.  A compact CY\(_3\) denominator theorem
requires further host data:
\[
  (X_5,\widetilde{\Gamma}_5,\mathcal L_5,Z_5),
\]
where \(X_5\) carries reduced \(\DT\) orientation theory,
\(\widetilde{\Gamma}_5\) makes the multiplier honest,
\(\mathcal L_5\) has weight \(c_5(0)/2\), \(Z_5\) is the reduced
\(\DT\) or chiral character in that line, and the BKM root
multiplicities come from the same Jacobi input.

\subsection*{Formal fourth roots}

The formal expression
\[
  \Phi_5^{1/2}:=(\Phi_5^{\mathrm{int}})^{1/4}
\]
has line-bundle weight \(2/4=1/2\).  An automorphic form of this weight
requires a cover \(\widetilde{\Gamma}\), an automorphy factor, a
character or multiplier \(\rho\), a branch convention, and divisor
multiplicities compatible with the chosen root:
\[
  \rho^4=\chi_{\Phi_5^{\mathrm{int}}},\qquad
  \operatorname{div}(\pi^*\Phi_5^{\mathrm{int}})
  \in 4\,\operatorname{Div}(\widetilde{\Gamma}\backslash\mathbb H_2).
\]
The metaplectic double cover \(\Mp_4\to\Sp_4\) supplies half-integral
weights.  A fourth-root character for \(\Phi_5^{\mathrm{int}}\) remains
additional data.  A weight-\(1/2\) BKM row therefore begins with a
separate Jacobi or vector-valued input of zeroth coefficient \(1\),
together with its multiplier, product chamber, and divisor theorem.

\subsection*{Order-$5$ host criterion}

The primitive CHL \(K3\times E\) denominator construction uses the
elliptic origin-fixing cyclotomic multiplier.  Its arithmetic constraint
is
\[
  \varphi(N)\mid 2,
\]
which gives \(N\in\{1,2,3,4,6\}\).  Since \(\varphi(5)=4\), the
order-$5$ automorphic row lies outside this primitive CHL construction.

A \(5\)-torsion translation on \(E\) preserves \(dz\) and can help make a
quotient free, but it supplies no origin-fixing cyclotomic multiplier
line.  A separately constructed quotient or Borcea--Voisin host may
supply the geometric component \(X_5\).  The denominator theorem still
requires the cover \(\widetilde{\Gamma}_5\), the automorphic line
\(\mathcal L_5\), the reduced \(\DT\) character \(Z_5\), and the
root-multiplicity comparison listed above.

Consequently the fourth-root identity
\[
  \bigl(Z^{X_5}_{\DT,\red}\bigr)^{-1/4}=\Phi_5^{1/2}
\]
becomes a theorem exactly after those four data are constructed in the
same normalisation.  The scalar form \(\Phi_5^{\mathrm{int}}\) and its
formal fourth root provide the automorphic target while leaving the
compact source and the partition-function equality as required data.

\subsection*{\(K3\times E\) invariants}

Whenever \(K3\times E\) appears in this comparison, the compact
invariants remain distinct:
\[
  \kcat(K3\times E)=0,\qquad
  \kch(K3\times E)=0,\qquad
  \kch^{\mathrm{Heis}}(K3\times E)=3,\qquad
  \kBKM(\mathfrak g_{\Delta_5})=5,\qquad
  \kfib(K3)=24.
\]
The values come from different constructions: K\"unneth on the compact
total space, compact Hodge supertrace, the Heisenberg--Mukai
specialisation, Borcherds weight, and Mukai-lattice rank.  The
order-$5$ boundary value \(\kBKM(\Phi_5^{\mathrm{int}})=2\) is an
automorphic boundary constant separate from the compact \(K3\times E\)
invariants.

\subsection*{Conclusion}

At order \(5\) the proved automorphic statement is
\[
  c_5(0)=4,\qquad
  \kBKM(\Phi_5^{\mathrm{int}})=2,\qquad
  \wt(\Phi_5^{\mathrm{int}})=2.
\]
The half-integral weight \(1/2\) belongs to a separate metaplectic or
multiplier-cover construction with its own input coefficient, branch,
and divisor theorem.  The primitive CHL table, the order-indexed
Nikulin/Gritsenko ladder, and the Gritsenko--Cl\'ery triple catalogue
therefore give three separate normalisations.  Mixing them changes the
claim.

\paragraph{Anchors.}
Nikulin 1980 and Mukai 1988 fix the symplectic K3 order list and
order-$5$ fixed-point data.  Eguchi--Hikami 2011 gives the twined K3
elliptic-genus table.  Gritsenko 1994, Gritsenko 1999, and Borcherds
1998 give the automorphic product weight formula.  Gritsenko--Cl\'ery
2008 gives the eight diagonal-divisor triples and the two
half-integral multiplier rows.  Oberdieck--Pandharipande 2018 and
Bryan--Oberdieck--Pandharipande--Yin 2018 cover the primitive CHL
\(K3\times E\) orders \(1,2,3,4,6\).
